\color{unimain}{\subsection{Логика отключения смежного трансформатора}}
\color{black}

\begin{enumerate}

\item
Ввод в работу логики отключения смежного трансформатора осуществляется с помощью программного переключателя \textbf{«ЛО см Т от ТЗНП»}. 

\item
Алгоритм работы логики отключения смежного трансформатора приведен на рисунке \ref{pict:LO_T}, уставки функции приведены в таблице \ref{app_set:losmt}.

\vspace{1mm}
\begin{figure}[H]
\centering
\includegraphics[width=1\textwidth,height=1\textheight,keepaspectratio]{\fbpath/07.01 РЗТ/021.0361 ЛО смеж Т/img/LO_T.pdf}
\captionof{figure}{Функциональная схема логики отключения смежного трансформатора}\label{pict:LO_T}
\end{figure}

\item
Оперативный ввод/вывод логики отключения смежного Т от ТЗНП производится с помощью БУ (блока управления функциями) \textit{«Вывод ЛО см. Т»}. 

\item
С целью исключения возникновения недопустимого режима работы параллельного трансформатора с изолированной нейтралью при замыкании на землю одной фазы. Предусматривается возможность первым действием сформировать сигнал \textit{«Откл. см. Т от ТЗНП»} на отключение выключателей стороны ВН смежного трансформатора. Логика отключения смежного трансформатора формируется по пуску ТЗНП ВН с независимой выдержкой времени, которая регулируется уставкой \textbf{«Тср откл см Т от ТЗНП»}.


\end{enumerate} 